% AI-written, AI cross-reviewed; no human sign-off; not a claim that RH is proved
% RH Notes blueprint source excerpts for Z1.3.03
% Verbatim excerpts in the notation of the parent blueprint. The labelled editions and source ranges appear below.
% See the unit YAML for the current proof status. Earlier-round passages retain their original wording.

% Start of atlas edition annotation. The following block is outside the original excerpt.
\paragraph{Status of named inputs in this reading.}
\noindent\nolinkurl{GlobalTheory.T.exists_admissible_hecke_projection_nonvanishing} --- in-repo constant (\path{Development/AbelianVarietyTheory/Proofs/FoundationProofs/ExistsAdmissibleHeckeProjectionNonvanishingStatement.lean:13--49}); not compared.\par
External theorem numbers in the original excerpt are not verified here unless a separate edition-specific source record documents the comparison.
% End of atlas edition annotation.

\section*{Z1.3.03. Equivariant Hecke projection}

% Parent source SHA-256: 1af64922b98ed2310744411b1f18b25cc426c97220646321471836e1a3a409d6
% blueprint/ch_P.tex lines 12473 through 12485; r5:small:Z1.3.03
\subsection{Z1.3.03. Integral Hecke projection and the functional}
\label{r5:small:Z1.3.03}
\textit{Partial.} The descent of the period continues on page~\pageref{r6:Z1.3.03}. Write $w=2e-1$, with $e\ge1$, and $B=X.B(N,\ell,w)$. The data include the ordinary normal form, holomorphic tower, $\ell\ne2$, $\ell\nmid N$, $N\ge5$, the bounded place embedding, and the local formulas for the two diagonal entries and the lower-left entry of $\rho$. The compatible maps $\pi_i,\sigma_i$ split the admitted Hecke polynomial $P$ on $V_B$, with $P(\mathrm{ev})=1$. The finitely supported functional $\lambda$ has inertia-fixed coefficients relative to $c$, and its level-$n$ Hecke congruence on the $\chi$-part loses at most a fixed $c_2$. The conclusion gives $a$, a common coefficient ring $c_f.O\supset c.O$ containing the matrix entries of $W$, and an additive $\mathbb Z_\ell$-linear map $p_r$. Its coordinates on $B$ lie in $c_f.O$. There it intertwines Galois, $T_q$ for $q\nmid\ell N$, the diamond operators for $(d,N)=1$, and $U_\ell$. Its second coordinate and $\ell^a\lambda$ satisfy the same period congruences in $\mathrm{IsLamFR2}$. The displayed proposition permits $\lambda=0$.
\par{\footnotesize \nolinkurl{GlobalTheory.T.exists_admissible_hecke_projection_nonvanishing} (\nolinkurl{Development/AbelianVarietyTheory/Proofs/FoundationProofs/ExistsAdmissibleHeckeProjectionNonvanishingStatement.lean}\allowbreak:13--49). \nolinkurl{GlobalTheory.KummerEisensteinGalois.IsLamFR2} (\nolinkurl{Development/AbelianVarietyTheory/Definitions/Foundations/IsLamFR2.lean}\allowbreak:13--16).\par}

On the diamond $\chi$-submodule, the commutation of Hecke and Galois operators with the diamonds preserves the character condition. Let $L$ be the limiting functional on the scalar extension, approximating the finite-level values with error $C|\ell|^n$. Subtract $\mathrm{ev}(T)L(b)$ from $L(Tb)$, insert the two approximants, and use the assumed congruence between them. All three errors tend to zero for each admitted $T$. The resulting eigenidentity persists on the scalar extension by linearity. Applying it successively to words, then to $P$, gives $L(Pb)=L(b)$. On each summand one uses $L(\sigma_i v)=u\kappa_i v_1$ in the source's zero-based coordinates. Thus $p_0(b)=\sum_i\kappa_i\pi_i(1\otimes b)$ satisfies $L(b)=u p_0(b)_1$ on the character part. The lower row $(0,1)$ of inertia fixes this coordinate. This allows comparison of the inertia-transformed finite-level approximant with the original one at every precision.
\par{\footnotesize \nolinkurl{GlobalTheory.LamChi.apply_poly_of_stable} (\nolinkurl{Development/AbelianVarietyTheory/Proofs/FoundationProofs/ApplyPolyOfStable.lean}\allowbreak:109--138). \nolinkurl{GlobalTheory.LamChi.exists_pr_isLamFR} (\nolinkurl{Development/AbelianVarietyTheory/Proofs/FoundationProofs/ApplyPolyOfStable.lean}\allowbreak:169--287). \nolinkurl{GlobalTheory.LamPr.exists_period} (\nolinkurl{Development/AbelianVarietyTheory/Proofs/FoundationProofs/ExistsRestrict.lean}\allowbreak:174--201).\par}

Here is the normalization step, including the zero case. If $L$ vanishes on the character part take zero multiplier and periods one. Otherwise choose $b_0$ with $y=L(b_0)\ne0$. From fixed algebraic approximants $z_n$ with $|z_n-y|\le |\ell|^{n+1}|y|$ we get $|z_n|=|y|$. The elements $z_n/z_0$ are fixed units converging to $y/z_0$. Multiplication of $p_0$ by $z_0/p_0(b_0)_1$ gives the period $y/z_0$ for all $b$. Consecutive period approximants agree to precision $n$ by the ultrametric inequality. For the resulting $p_0$ (absorbing that multiplier), choose a finite free basis of $B$. Choose $a$ so that $\ell^a$ as a positive real number dominates the sum of the norms of the two coordinates of $p_0$ on every basis vector. Since $|\ell^a|_\ell=\ell^{-a}$, all the scaled coordinates are integral. Adjoin these finitely many elements to one integral coefficient ring containing $c.O$ and the matrix entries. Linearity gives integrality on $B$. Scaling both $p_0$ and $\lambda$ preserves the period bound since $|\ell^a|_\ell\le1$. The equivariance of each $\pi_i$ gives that of the sum, including the separate $U_\ell$ identity.
\par{\footnotesize \nolinkurl{GlobalTheory.LamChi.exists_fixed_period} (\nolinkurl{Development/AbelianVarietyTheory/Proofs/FoundationProofs/ApplyPolyOfStable.lean}\allowbreak:76--99). \nolinkurl{GlobalTheory.LamChi.exists_pr_isLamFR} (\nolinkurl{Development/AbelianVarietyTheory/Proofs/FoundationProofs/ApplyPolyOfStable.lean}\allowbreak:289--333). \nolinkurl{GlobalTheory.AdmProj.exists_admissible_hecke_projection_nonvanishing_alt} (\nolinkurl{Development/AbelianVarietyTheory/Proofs/FoundationProofs/ExistsAdmissibleHeckeProjectionNonvanishing.lean}\allowbreak:74--183).\par}

\textit{[Proved in Lean, not reproduced:]} the modular-tower construction of $L$, with one unit $u$ and algebraic $\kappa_i$ on all splittings, and the fixed approximation theorem with bounded precision loss. The latter says that an integral algebraic $x$ fixed modulo $\ell^k$ by the specified inertia subgroup is congruent modulo $\ell^{k-c_1}$ to an integral exactly fixed $x'$, with $c_1$ depending only on the starting coefficient ring. These inputs are used before the finite-basis scaling above. \textit{[Hypothesis, discharged at the use:]} $P,\pi_i,\sigma_i$ and the congruence are supplied as arguments of the projection statement. The global Kato construction uses this very statement as its third input.
\par{\footnotesize \nolinkurl{GlobalTheory.LamPr.exists_period} (\nolinkurl{Development/AbelianVarietyTheory/Proofs/FoundationProofs/ExistsRestrict.lean}\allowbreak:174--188). \nolinkurl{GlobalTheory.T.exists_inertiaFixed_near_almostFixed_coefficients} (\nolinkurl{Development/AbelianVarietyTheory/Proofs/FoundationProofs/ExistsInertiaFixedNearAlmostFixedCoefficientsStatement.lean}\allowbreak:6--9). \nolinkurl{UpperTheory.V.Z.Kato.Final.zetaUr} (\nolinkurl{Development/AutomorphicArithmetic/AutomorphicVarieties/ModularityArgument/KatoQuadraticCharacters/UnramifiedZeta/ZetaUr.lean}\allowbreak:9--12).\par}

% blueprint/ch_P.tex lines 12977 through 12991; r6:Z1.3.03
\subsection{Z1.3.03. Descent of the common period and uniform loss}\label{r6:Z1.3.03}
Continue the integral projection argument of Section~\ref{r5:small:Z1.3.03}. For $b$ in the integral Betti lattice put $a_n(b)=\lambda((j(\phi_n[i]))_i)$, using a cusp representative $\phi_n$ at precision $n$ and denominator $D=1$. Because $\lambda$ uses only finitely many coefficients there is a bound $C$ on the unit ball. Consecutive representatives have coefficient differences bounded by $|\ell|^n$. Thus, for $t>n$, $|a_t(b)-a_n(b)|\le C|\ell|^n$, independently of the choices. Taking limits in $\mathbb C_\ell$ gives an additive $\mathbb Z_\ell$-linear functional on the lattice, by the finite-level addition formula for representatives. It extends to the scalar span and then linearly to the ambient tensor space. The bound for every representative survives the limit. If $\tau$ fixes the finitely many coefficients of $\lambda$, restrict it along $j$ to obtain $s$; the cusp representative of $sb$ is the coefficientwise translate of that of $b$. Their two approximations have error tending to zero, hence $L(sb)=\tau L(b)$ on the lattice.
\par{\footnotesize \nolinkurl{GlobalTheory.LamPr.exists_limit} (\nolinkurl{Development/AbelianVarietyTheory/Proofs/FoundationProofs/ExistsRestrict.lean}\allowbreak:108--172). \nolinkurl{GlobalTheory.LamPr.exists_restrict} (\nolinkurl{Development/AbelianVarietyTheory/Proofs/FoundationProofs/ExistsRestrict.lean}\allowbreak:69--75).\par}

For $w=2e-1$, the vectors $\sigma_i(1,0)$ are eigenvectors for the character $\epsilon_\ell^{w+1}$ of the inertia subgroup fixing the prescribed finite set. To descend the eigencondition, intersect the kernels of $g_s-\epsilon_\ell(s)^{w+1}$ on the rational span of the finite free lattice. A finite intersection already has the least possible dimension. Base change preserves this kernel, and multiplying by powers of $\ell$ clears the denominators of a rational basis. It suffices to kill the resulting integral eigenvectors. On such a vector compare the translate of a cusp representative with its scalar multiple by the character. Uniqueness modulo $\ell^n$ gives the corresponding transformation of every coefficient $a$; the integral Tate inertia bound gives $\ell^c a\equiv0\pmod{\ell^n}$ with $c$ independent of $a,n$. The coefficient and approximation bounds now give $|L(b)|\le C|\ell|^{n-c}$ for all $n\ge c$, hence $L(\sigma_i(1,0))=0$.
\par{\footnotesize \nolinkurl{GlobalTheory.LamPr.exists_period} (\nolinkurl{Development/AbelianVarietyTheory/Proofs/FoundationProofs/ExistsRestrict.lean}\allowbreak:174--230). \nolinkurl{GlobalTheory.ThKill.discrete} (\nolinkurl{Development/AbelianVarietyTheory/Proofs/FoundationProofs/ThetaKillsFplusDiscreteForm.lean}\allowbreak:83--133). \nolinkurl{GlobalTheory.T.integral_tate_inertia_cyclotomic_twist_torsion} (\nolinkurl{Development/AbelianVarietyTheory/Proofs/FoundationProofs/IntegralTateInertiaCyclotomicTwistTorsionStatement.lean}\allowbreak:7--11).\par}

Write $y_i=L(\sigma_i(0,1))$. When all $y_i=0$ take $u=1$, $\kappa_i=0$. Otherwise fix $y_j\ne0$. Choose one finite extension $K/\mathbb Q_\ell$ containing the coefficients of $\lambda$ and all coefficients of tensor expansions of the finitely many $\sigma_i(0,1)$ in lattice elements. Any $\tau\in G_K$ restricts along $j$ to an $s$ as above; thus $\tau y_i=\rho(s)_{11}y_i$ by the first-coordinate vanishing, with the same multiplier for every $i$. It is nonzero since $\tau y_j\ne0$. The ratios $y_i/y_j$ are fixed; Ax's fixed-field theorem puts them in the complete finite field $K$. By density, choose algebraic $a$ with $|a-y_j|<|y_j|$. Then $|a|=|y_j|$. Set $u=y_j/a$ and $\kappa_i=a(y_i/y_j)$. Linearity gives $L(\sigma_i v)=u\kappa_i v_1$ simultaneously and $|u|=1$.
\par{\footnotesize \nolinkurl{GlobalTheory.line_functional_eq_period_mul} (\nolinkurl{Development/AbelianVarietyTheory/Proofs/FoundationProofs/LineFunctionalEqPeriodMul.lean}\allowbreak:23--99). \nolinkurl{GroundTheory.PAdicHodge.exists_eq_coe_of_forall_smul_eq} (\nolinkurl{GroundTheory/PAdicHodge/AxLemma.lean}\allowbreak:530--537).\par}

For the second loss, let $H$ consist of automorphisms fixing $c.O$ pointwise and inducing the identity on the residue field. If an integral algebraic $x$ satisfies $|s x-x|\le\ell^{-k}$ for $s\in H$, Ax's uniform orbit bound gives an algebraic $y$ fixed by $H$ with $|x-y|\le C_\ell\ell^{-k}$. To apply the field form to $H$, take its closure: $x$ has finite orbit, so the same bound holds on that closure. The algebraic fixed field is henselian over $\mathbb Q_\ell$ with residue characteristic $\ell$. Choose $c_1$ so that $C_\ell<\ell^{c_1}$. For $k>c_1$ the estimate implies $|y|\le1$ and $x\equiv y\pmod{\ell^{k-c_1}}$; for $k\le c_1$ take $y=0$ (truncated subtraction). This closes the fixed-approximation step, uniformly in $k,x$. The normalization and denominator clearing of Section~\ref{r5:small:Z1.3.03} now apply. The published estimates and the remaining input are recorded in Section~\ref{r6:kato-lit}.
\par{\footnotesize \nolinkurl{GlobalTheory.PerFix.inertiaC} (\nolinkurl{Development/AbelianVarietyTheory/Proofs/FoundationProofs/ExistsInertiaFixedNearAlmostFixedCoefficients.lean}\allowbreak:16--37). \nolinkurl{GlobalTheory.PerFix.exists_fixed_near} (\nolinkurl{Development/AbelianVarietyTheory/Proofs/FoundationProofs/ExistsInertiaFixedNearAlmostFixedCoefficients.lean}\allowbreak:39--71). \nolinkurl{GlobalTheory.AxLemma.ax_lemma} (\nolinkurl{Development/AbelianVarietyTheory/Proofs/IwasawaDecomposition/BoundedOrbitHasNearbyFixedPoint.lean}\allowbreak:211--221).\par}

\emph{Partial. [Proved in Lean, not reproduced:]} on the specific modular tower, for every $b\in B$, $n\ge0$ and odd $w$, construct a cusp representative satisfying both the level-$n$ comparison and integral exponents, and prove its coefficientwise uniqueness at the needed precision. The Tate input used above is the uniform integral estimate for all $x\in A$ with $s x\equiv\epsilon_{\ell,n}(s)^{w+1}x\pmod{\ell^n}$ on inertia fixing $F$; the bound $c$ may depend on $F,w$, but not on $n,x$. These two concrete upstream comparisons remain open in this exposition. The ordinary input, polynomial splitting, and congruences for the specified $\lambda$ are hypotheses of the unit.
\par{\footnotesize \nolinkurl{GlobalTheory.T.exists_cuspFunctional_oddWeight_oddPrime} (\nolinkurl{Development/AbelianVarietyTheory/Proofs/FoundationProofs/ExistsCuspFunctionalOddWeightOddPrimeStatement.lean}\allowbreak:8--9). \nolinkurl{GlobalTheory.LamPr.exists_limit} (\nolinkurl{Development/AbelianVarietyTheory/Proofs/FoundationProofs/ExistsRestrict.lean}\allowbreak:119--139). \nolinkurl{GlobalTheory.ThKill.discrete} (\nolinkurl{Development/AbelianVarietyTheory/Proofs/FoundationProofs/ThetaKillsFplusDiscreteForm.lean}\allowbreak:95--105). \nolinkurl{GlobalTheory.T.integral_tate_inertia_cyclotomic_twist_torsion} (\nolinkurl{Development/AbelianVarietyTheory/Proofs/FoundationProofs/IntegralTateInertiaCyclotomicTwistTorsionStatement.lean}\allowbreak:7--11).\par}

% Parent source SHA-256: 1fdc6ea9aac6cf10316f85b3253b33ae98289a76ff5aa2e79daaaf11e057ce18
% blueprint/ch_P_r7.tex lines 666 through 750; r7:Z1.3.03
\subsection{Z1.3.03. Uniform integral Tate loss and the cusp}\label{r7:Z1.3.03}
\textbf{A01, published vanishing and its integral consequence.}
Use J.~Tate, \emph{\(p\)-Divisible Groups}, in \emph{Proceedings of a
Conference on Local Fields} (1967), \cite[\S3.3, Theorem~2, p.~176]{Tate-pdiv}. The result used is
vanishing of \emph{both} continuous \(H^0\) and \(H^1\) for a
one-dimensional twist of the completed algebraic closure whose character
cuts out, after finite base extension, a totally ramified
\(\mathbb Z_\ell\)-extension.
\href{https://www.math.purdue.edu/~tongliu/teaching/598/p-divisible.pdf}{Tate, original text.}
The standalone Lean lemma covers every prime \(\ell\). Its Kato application assumes \(\ell\) odd. Fix \(A\) with \(v_A(\ell)>0\), a positive twist \(j\) and the finite set \(F\). Embed
\(\Omega\) at \(A\) in \(\overline{\mathbb Q}_\ell\). Concretely, on each finite \(E/\mathbf Q\) in \(\Omega\), complete at the prime induced by \(A\) and embed the completion in \(\overline{\mathbf Q}_\ell\). The resulting finite nonempty sets of compatible-place embeddings have surjective restriction maps, since an embedding of a finite local field extends across a finite extension. Compactness of their inverse limit gives one embedding with inverse image of the unit ball exactly \(A\). See \nolinkurl{ReciprocityTheory.Neutral.Cmp.exists_compat}, \path{LocalFieldAutomorphisms/ValuationIntCastLtOneIff.lean:83--89}. Inside
\(\mathbb C_\ell\), take \(K\) to be the completion of the maximal
unramified extension of \(\mathbb Q_\ell\), with the images of \(F\) and
\(\mu_\ell\) adjoined (adjoin \(\mu_4\) also for \(\ell=2\)).
This finite extension of a complete discretely valued field is complete
and discretely valued, of mixed characteristic and with algebraically
closed residue field. The algebraic closure of \(K\) inside
\(\mathbb C_\ell\) contains \(\overline{\mathbb Q}_\ell\), so its
completion is \(\mathbb C_\ell\).
Each \(s\in G_K\) preserves the embedded \(\Omega\), the subfield
of elements algebraic over \(\mathbb Q\). It restricts to an inertia
element fixing \(F\): the images of \(F\) are in \(K\), and the maximal local unramified extension it fixes has residue field \(\overline{\mathbf F}_\ell\). Thus \(s\) fixes the reduction of every integral element. It preserves the valuation and fixes all residues on restriction, which is exactly inertia.
See \nolinkurl{ReciprocityTheory.Neutral.Cmp.res_mem_ine_iff}, \path{LocalFieldAutomorphisms/Inertia.lean:45--50}. Thus the source hypothesis bounds every such \(s\).
\par The cyclotomic image lies in principal units (in \(1+4\mathbb Z_2\)
at 2), where the logarithm identifies closed infinite subgroups
with copies of \(\mathbb Z_\ell\). The image is infinite: otherwise one finite
extension of \(K\) would contain every \(\ell\)-power root of unity,
contradicting discreteness since
\(v(\zeta_{\ell^n}-1)=1/\varphi(\ell^n)\) by Eisenstein. Its
\((-j)\)-th power still has infinite image: under the logarithm multiplication by the nonzero integer \(-j\) is injective on this torsion-free group.
Every finite layer is totally ramified since the base is complete
discretely valued with algebraically closed residue field. These checks
give the hypotheses of Theorem~2 for \(V=\mathbb C_\ell(\chi^{-j})\).
\par Put \(Dx(s)=\chi(s)^{-j}s(x)-x\). By the two vanishings,
\(D:V\to Z^1_{\rm cts}(G_K,V)\) is a continuous bijection of
\(\mathbb Q_\ell\)-Banach spaces: continuous functions on the compact
group have the supremum norm and the cocycle identities are closed.
The dilates by \(\ell^{-n}\) of the image of the closed unit ball cover
the target. Baire gives a ball about zero in the closure of the image
of the unit ball, since this closure is an additive subgroup. For a
target in that ball approximate by an image to error in \(\ell\)
times the same ball. Divide the error by \(\ell\), repeat, and add
the successive preimages multiplied by \(1,\ell,\ell^2,\ldots\).
Completeness gives a preimage in the unit ball. Thus
\(|x|\le\ell^c\sup_s |Dx(s)|\) for a single nonnegative \(c\).
For integral \(x\), replacing \(\chi(s)\) by its residue representative
modulo \(\ell^n\) changes the untwisted product by at most \(\ell^{-n}\)
for \(n>0\), and \(n=0\) follows from integrality. The assumed congruence on inertia fixing \(F\) therefore bounds \(\sup_{G_K}|Dx|\) by \(\ell^{-n}\).
Thus \(|\ell^c x|\le \ell^{-n}\) for every such \(x,n\) with the same \(c=c(\ell,A,j,F)\). Dividing by \(\ell^n\) gives an element of \(A\), exactly the integral congruence.
\nolinkurl{GlobalTheory.T.integral_tate_inertia_cyclotomic_twist_torsion},
\path{FoundationProofs/IntegralTateInertiaCyclotomicTwistTorsionStatement.lean:7--11};
\nolinkurl{GlobalTheory.IntTate.integral_tate_inertia},
\path{FoundationProofs/IntegralTateInertiaCyclotomicTwistTorsion.lean:524--603}.

\textbf{C02, hand proof.} The Hilbert~90 step in descending the cusp
lift allows a finite-index subgroup \(\Delta\) which may be nonnormal.
Let \(L\) be the subfield fixed by \(\Delta\) and \(D\) its pointwise
stabilizer. Since \(\Delta\subseteq D\), the index of \(D\) is finite. The cocycle \(c:\Gamma\to K^\times\) for the given action on the ambient field \(K\) equals one on \(D\)
by hypothesis. Choose left-coset representatives for \(\Gamma/D\).
Their restrictions to \(L\) are distinct field embeddings. To see
independence, take a shortest nonzero linear relation. Applying it at
\(bx\) and subtracting the relation at \(x\) times the first embedding
of \(b\) removes one coefficient; choose \(b\) separating the first and
another embedding. This contradicts minimality (a single embedding
is nonzero at 1). Thus for some \(x\in L\) the sum
\(y=\sum_t c(t)t(x)\) is nonzero. The summands depend only on
cosets since \(D\) fixes \(x\) and \(c(td)=t(c(d))c(t)=c(t)\) for every \(d\in D\).
Left multiplication permutes the cosets, so \(c(s)s(y)=y\).
\nolinkurl{GlobalTheory.ThLift.hilbert_ninety},
\path{FoundationProofs/ExistsInvariantLift.lean:13--55}.

\textbf{B04.} For every admitted odd-weight \(b\), construct a cusp
representative at precision \(n\) with its prescribed class, invariance,
complementary power condition and integral Fourier lattice.
\textbf{B05.} For two representatives satisfying the successive-level
hypotheses, prove the coefficient congruence at the stipulated precision.
The source uses existence and uniqueness separately. The published
period vector in Kato 12.5(1) belongs to a modular realization. It
provides neither a representative in the supplied tower nor a uniqueness
estimate for its Hahn coefficients. C02 accounts for the Hilbert~90
step in B04.
\nolinkurl{GlobalTheory.ThEx.thex},
\path{FoundationProofs/ExistsCuspFunctionalOddWeightOddPrime.lean:16--30};
\nolinkurl{GlobalTheory.T.cgA_coeff_of_isThB_succ},
\path{FoundationProofs/IsThBSuccCoeffCongStatement.lean:8--10}.

